\section{Cross-geometry and photo-patch diagnostics}
\begin{figure*}[t]
\centering
\includegraphics[width=\textwidth]{../../experiments/results/geometry_prediction_wide/geometry_prediction.pdf}
\caption{Four tests on identical geometry coordinates. Dark means the
test holds. The first two panels use 1200 logarithmically sampled
variances between $v_L$ and 100, plus the analytic infinite-variance
endpoint for the first panel. They concern the fixed class-1 corner
$(.3,-.01)$ and noiseless data. The third panel checks the analytic
class-1 channel conditions for $[.69,.71]$ and 25 calls, with
$\delta=10^{-4}$. Light means unclassified by those sufficient
conditions. The fourth panel checks exact 106-call iterates for all
three noiseless template observations, with radius-$.1$ success.
All panels fix $\alpha=.001,a=1,r=1.29$. The grid is logarithmic in
$h$ and linear in $d$. It was widened after an initial narrower grid
showed no negative points in the first two tests. The grid defines
a diagnostic sample and supplies no continuous-parameter guarantee.}
\end{figure*}
The four masks have respectively 995, 1673, 1383 and 1248 positive
points out of 1681. Every channel-certified point enters the class-1
prepared rectangle. Only 1176 of these points recover all three
classes. Other classes' entry, terminal distance gaps and available
refinement budget require separate tests as specified in
Theorem~\ref{thm:criterion}.

\subsection{Matched risk illustrations}
\begin{table}[t]
\centering\small
\setlength{\tabcolsep}{3pt}
\begin{tabular}{lrrrr}
\toprule
Model & $K$ & Preset & Restricted & Terminal\\
\midrule
Calibrated & 106 & $1.53\cdot10^{-12}$ & $.3163$ & $.3163$\\
Separate noise & 58 & $1.28\cdot10^{-8}$ & $.3169$ & $.3169$\\
Fine calibration & 154 & $1.47\cdot10^{-20}$ & $.3169$ & $.3169$\\
\bottomrule
\end{tabular}
\caption{Prior-weighted squared error of the raw RED iterates on paired draws.
All methods within a row share the observation, initialization, total
budget and terminal variance.}
\label{tab:risk_diagnostics}
\end{table}
We fix $a=1$, $d=2$, $r=1.29$ and $h=1.78$ in all three rows of
Table~\ref{tab:risk_diagnostics}. The calibrated row uses
$\alpha=10^{-3}$, $\tau=\nu=\alpha^2$ and $\theta=0$.
The second row uses $\alpha=10^{-6}$, $\tau=10^{-4}$,
$\nu=10^{-5}$, $v_L=(\tau^2+\nu^2)^{1/2}$ and $\lambda=1$.
The preset uses $v=.70$ for its observation-selected intermediate calls.
The restricted policy alternates $.17$ and $2.5$ for 25 calls, then
uses $v_L$. The terminal comparator uses $v_L$ throughout.
The fine row repeats the calibration at $\alpha=10^{-5}$.
Budgets follow the respective refinement formulas.

We draw 512 independent pairs per class per row using seed 20261109
and weight class averages by $\pi$. All settings are prespecified.
All draws satisfy the coordinate bounds in~\eqref{eq:ind_event}.
In all three rows, mode accuracy and radius-$.1$ template-neighborhood success
are both $1$ for the preset and $.683$ for the two comparators.
We compute these metrics separately from squared error.
The table illustrates the separation.
The guarantees for the continuum of variances and arbitrary restricted
switching follow from the common barrier proof.


\paragraph{Fine excess risk.}
For the fine row, the analytic oracle risk is $1.49999\cdot10^{-20}$.
We estimate the excess risk directly as the prior-weighted mean of
$\norm{\widehat x-m_Z(1)}^2$, which is $8.49\cdot10^{-32}$ in double precision. This is an implementation
reading near numerical resolution, dominated by rounding. Its specific
value does not measure agreement with the $5\cdot10^{-30}$ bound.
Repeating the intermediate steps at 50 and 80 digits gives
$3.48001\cdot10^{-34}$ at both precisions using exact decimal parameters
and the same binary64 normal draws. After verifying terminal entry for
every sample, we propagate the affine residual
$[I-(B+I)/4]^n(x_h-m_Z(1))$.
The omitted terminal leakage changes the output by at most
$12e^{-.03/v_L}<10^{-130288343}$.
Here $m_Z(1)=(B+I)^{-1}(b+\mu_Z)$; labels enter evaluation only. The ordinary empirical MSE fluctuates on a
larger scale than this excess.

\paragraph{Precision reference.}
The program \path{experiments/code/precision_diagnostic.py} replays
the saved normal draws at 50 and 80 decimal digits. Parameters are
exact decimal constants; the standard-normal draws are the same
binary64 numbers in the original sampling stream. For each of 1536
samples it checks terminal entry before using the analytic affine
recursion. The terminal error bound $12\exp(-.03/v_L)$ comes from
the uniform distance gap $.06$ and rectangle invariance. Its
base-10 logarithm is less than $-130288343$ here. Both precisions
give excess risk $3.48001105532541\cdot10^{-34}$.
This controls the omitted soft-template term but does not constitute
an interval enclosure of the arbitrary-precision arithmetic.

\paragraph{Initialization and weak-coefficient boundaries.}
The least-squares start lies inside the correct terminal rectangle
when $|e_1|\le .1\alpha$ and $|e_2|\le.005$.
For exact updates with $\lambda=\alpha/4$, class-1 drift at $s=.3$
is strictly outward on $E$ for every admissible scale, since
$b_1-\alpha s\ge .45\alpha$ and
$c(v)(m_v-s)\ge-(1.02+.3)\alpha/4$.
Changing the data loss to a likelihood-normalized form would also
change the drift and step-size conditions.

\paragraph{Weak-regularization budget.}
For noiseless observations of the three nominal templates,
$\alpha=10^{-5}$ and $\lambda=\alpha/4$, terminal-only iteration
from the backprojection has prior-weighted squared error $.75362$,
$.74998$, $.71238$, $.42955$ and $.001932$ at budgets
154, $10^3$, $10^4$, $10^5$ and $10^6$.
The first radius-$.1$ entry times are 1, 828881 and 736823.
These are observed trajectories at a changed coefficient, not
uniform hitting-time guarantees. See
\path{experiments/code/weak_budget_diagnostic.py}.

\paragraph{Patch data and protocol.}
The source images are the two photos distributed in
\texttt{sklearn.datasets.load\_sample\_images}.
They are converted to grayscale and divided into nonoverlapping
$4\times4$ patches. A seed-20261110 subset of 2048 patches from
\texttt{china.jpg} trains eight means with $k$-means (five starts);
the masses are cluster frequencies and the common isotropic variance
is the mean squared training residual per coordinate.
A disjoint-photo subset of 256 patches from \texttt{flower.jpg}
is evaluated. No labels are defined for these empirical truths.
The exact MMSE denoiser of the fitted mixture uses total variance
$\tau^2+\sigma^2$ and the RED residual is unnormalized.
The blur applies periodic convolution by $[.25,.5,.25]$ on each
patch axis. This is a patch-level degradation operator.
All policies share the resulting noisy observations, backprojection,
154 calls, $\lambda=.1$, $\gamma=.25$ and terminal $\sigma=.02$.
Their mean squared Euclidean errors are respectively $.00990445$,
$.00988671$ and $.00986463$ for terminal, intermediate, and low/high
protocols. Per-pixel MSE divides these by 16. These are descriptive
results on one image pair with no schedule search or fitted-model
certificate. See \path{experiments/code/patch_diagnostic.py}.

Reproduction uses the repository's Python environment and a fresh
output directory for each of
\path{experiments/code/geometry_prediction_diagnostic.py},
\path{experiments/code/precision_diagnostic.py}, and
\path{experiments/code/patch_diagnostic.py}, each with
\texttt{--out <new-directory>}. Dependencies are NumPy, SciPy,
Matplotlib, Pillow, scikit-learn and mpmath.
